\chapter{개선 사항 및 제언}

\section{개선 사항}

\subsection{종합}

\subsubsection{일정 관련}

학기 말에 업무가 과도하게 집중되는 현상을 방지하기 위해 실질적인 조사 준비 과정을 조기에 착수해야 한다는 의견이 많았다. 이론 수업과 조사 준비를 병행하거나, 중간고사를 앞당겨 실시함으로써 본 조사 시기를 앞당기는 방안을 고려할 필요가 있다. 이를 통해 전사 및 보고서 작성에 필요한 절대적인 시간을 확보하고, 기말고사 기간과 과제 마감이 겹치는 부담을 완화하는 것이 나을 것이다. 

\subsubsection{총괄 관리자 선정}

수강 인원의 규모와 관계없이 프로젝트 전반의 역할을 분담하고 진행 상황을 관리할 총괄 관리자가 필요하다. 학기 초반부터 명확한 체계를 확립해야 조별 소통 혼선을 막을 수 있다. 이번 학기의 경우 총괄의 부재로 인해 전체적인 조율에 어려움이 있었다.

\subsubsection{사전 교육 구체화 및 예비 과제 도입}

모의 조사가 조교의 즉흥적인 대응에 의존하지 않도록 구체적인 상황 설정과 매뉴얼을 만드는 것이 좋을 것이라는 의견이 있었다. 또한, 모의 조사 영상을 함께 검토하며 객관적인 피드백을 주고받는 세션을 도입하면 좋을 것이다. 예비 과제를 늘려, 수강생들이 전사 과정의 어려움과 협업의 중요성을 사전에 충분히 체감할 수 있게 하는 것이 바람직하다.

\subsubsection{현장 조사 관련}

장거리 이동과 조사의 피로도를 고려하여 넉넉한 공간을 갖춘 차량을 대여해야 한다. 사전 답사 시에는 수강생이 최소 1인 이상 동행하여 현장의 분위기를 파악하는 것이 좋을 것이다. 자료제공인은 현지 섭외의 장점을 살리되, 불확실성을 최소화하기 위해 일부 인원은 사전에 섭외하는 것이 필요하다. 또한 현장에서의 돌발 상황에 대비해 녹음 장비는 최대한 여유 있게 준비해야 한다.

\subsection{문준모}

본 수업에서 가장 중요한 부분은 수강생들이 직접 현장에서 수행하는 언어 조사, 그리고 그에 뒤따르는 전사 및 정리 작업일 것입니다. 수강생들은 현장으로 향하기 이전까지 기록 언어학의 취지와 방법론에 대해 배운 뒤 어떻게든 썩 괜찮은 조사를 수행합니다. 그런데 정작 그 이후에 보고서를 만들기까지의 과정이 얼마나 손이 많이 가는지 미리 파악하기란 (선생님의 지속된 경고에도 불구하고!) 쉽지 않습니다. 과제의 볼륨을 늘리는 방향으로라도 자신의 전사 속도는 어느 정도인지, 보고서를 쓰는 데 정확히 어떤 요소가 필요한지, 팀원과 전사 내용을 맞춰 가는 감각은 어떤지를 체득할 기회를 얻었더라면 작업에 임하는 저의 태도가 어떻게 달라졌을지 하는 고민이 있습니다.

\subsection{박준영}

총괄이 있었다면 보다 원활하게 프로젝트를 수행할 수 있었으리라 생각합니다. 이번 언어조사에서 겪었던 가장 큰 어려움은, 필요한 역할을 분담하고, 각 조와 팀의 권한을 조정하고, 프로젝트 전반의 진행 경과를 확인하는 일이었습니다. 이 문제를 해결하기 위해서는 프로젝트 전반을 이끄는 총괄이 반드시 필요합니다. 2024년 언어조사에는 총괄에 해당하는 `반장'이 있어, 수강생 약 20명에 달하는 대형 언어조사를 성공적으로 이끌었던 것으로 보입니다. 올해는 조사자가 6명뿐이라 이런 `반장' 역할이 필요하지 않겠다는 다소 안일한 생각을 학기 초반에 했었는데, 언어조사를 준비하고 실행하는 과정에서 총괄의 빈 자리가 상당히 크게 느껴졌습니다. 이에, 다음해부터는 `언어조사 및 분석' 강좌의 규모에 관계없이 프로젝트 관리자를 한 명 두는 것을 제안합니다.

\subsection{유승민}

언어조사 과정 중 이용한 차량이 다소 협소했던 감이 있습니다. 이후의 언어조사에서는 조금 더 넉넉한 공간이 있는 차량을 대여한다면 편리할 것 같습니다. 

사전 조사에 최소한 수강생 1인이 참여하는 것이 좋을 것 같습니다. 이는 언어조사 당일에 겪을 혼란을 제하기 위함이며, 조교의 부담을 줄이는 것의 일환입니다. 또, 일정 구상 등이 더욱 이른 시기부터 이루어진다면 더욱 언어조사가 체계적으로 이루어지는 데에 도움이 될 것 같습니다.

자료제공인을 사전에 구하는 것이 단점이 많겠다는 생각이 들고, 이번 언어조사에서 자료제공인을 현지에서 수강생들이 직접 구한 데에 따른 이점이 많았다고 생각합니다. 다만 혹시 현지에서 자료제공인을 구할 수 있는 여건이 마련되지 않을 수도 있으니 자료제공인 중 일부는 사전에 섭외해놓는 것이 좋을 것 같습니다. 그리고 조사가 이루어지는 장소에 대해 더욱 체계적인 사전조사가 이루어졌으면 하는 생각이 있습니다.

\subsection{이상인}

이것이 제가 감히 올릴 말씀인지 잘 모르겠지만, 모의조사가 너무 조교의 즉흥에 크게 의존하는 것 같습니다. 나름의 매뉴얼과 설정을 잘 준비했겠지만, 강민하 조교님이 아니었더라면 어떻게 됐을지 모를 것 같기도 하네요… 여기에 더해 둘째로, 모의조사를 단순히 다른 팀이 관전하고 피드백하는 것을 넘어서서 함께 모여서 영상을 돌려 보는 세션이 있다면 좋을 것 같습니다. 조사를 직접 하고 있는 사람은 무엇을 잘했고 잘못했는지 스스로 되돌아볼 기회가 얼마 없을 것 같아서….

\subsection{이채현}

언어조사의 실질적인 준비과정이 보다 이른 시기부터 시작되는 게 좋을 거같습니다. 앞에 이론 부분을 집약적으로 진행함과 동시에 처음부터 언어조사를 단계적으로 진행해나가면서 이론 수업을 듣는 것이 더 효율적일 것 같습니다. 이론을 대부분 진행하고 실제 조사를 위한 작업을 착수하다보니까 점점 일정이 밀리는 것 같습니다. 또한 처음에 사전 조사와 질문지가 조금 더 이른 시기에 마감 기간을 두고 제작되어야 한다고 생각합니다.

중간고사를 먼저 실시하고 실제 언어조사를 좀 더 이른 시기에 다녀와서, 조사한 내용에 대해 더 여유있게 검토하고 전사를 하는 시간이 필요하다고 생각합니다. 조사를 10월 후반에 다녀오자마자 시험을 봐야하는 상황이었고, 수업 수강생의 인원이 적다보니 한 사람이 부담해야 할 양이 가중된 것으로 보입니다. 생각보다 한 번 전사하는 데도 오래 걸리고, 3번 정도 로테이션을 돌려야 하다 보니까 현실적으로 스케줄을 조율하기가 어려워 며칠밤을 새야 하는 상황도 많았습니다. 

\subsection{전지민}

팀별 인원수가 평상시에 비해 부족해짐으로 인해 단점이 많았다고 생각합니다. 수강생 수가 적어 어쩔 수 없는 일이지만, 1인당 전사 분량이 늘어난 것도 있고, 본 조사시에도 세 명이 어르신 세 명을 한번에 조사하려니 힘들었던 것 같습니다. 

전사 기한이 너무 촉박한 것 같다는 생각이 들었습니다. 특히 기말고사 기간과 겹치기도 하고, 저희는 기존에 비해 팀별 인원이 줄면서 더욱 촉박했던 것 같습니다. 한 주라도 시간이 더 있다면 좋지 않을까 생각합니다.

\section{다음 수강생을 위한 제언}

\subsection{종합}

\subsubsection{수강 신청 시 고려사항}

본 수업은 단순한 답사가 아닌 고강도의 학술 활동임을 인지해야 한다. 전공 학습량이 많거나 개인 일정이 바쁜 학기에는 수강을 지양하는 것이 좋으며, 수강 시에는 충분한 시간을 확보해야 한다. 사라져 가는 방언을 기록한다는 연구자로서의 사명감과 책임감이 필수적이며, 과정이 고되더라도 살아있는 언어를 다루는 값진 경험이라는 점을 기억해야 한다.

\subsubsection{빠른 작업 착수}
조사 종료 직후 즉시 전사 작업에 착수해야 후반부의 업무 과중을 피할 수 있다. 기말고사 기간 타 과목 기말고사로 인한 부담을 줄이기 위해서는 학기 중반까지 최대한 많은 분량을 처리해두는 전략이 필요하다. 계획 수립 시에는 `3시간 단위 목표 설정'이나 `6시간 간격 보고'와 같이 매우 구체적이고 강제성 있는 일정을 마련하여 작업이 지연되지 않도록 관리해야 한다.

\subsubsection{팀워크 중시와 철저한 장비 대비}

보고서 작성은 긴밀한 협력이 요구되는 거대한 팀 프로젝트이다. 개인의 태도가 팀 전체의 성과에 직결됨을 명심하고, 서로 격려하며 멘탈을 관리하는 자세가 필요하다. 현장 조사에서는 다수의 화자가 동시에 발화하거나 기기에 문제 발생하는 등 다양한 변수가 생길 수 있으므로, 녹음 장비는 가능한 한 넉넉하게 준비하여 만일의 사태에 대비해야 한다.

\subsection{문준모}

본 수업에서 무엇을 얻어 갈지는 상상 이상으로 본인 하기 나름이라고 느꼈습니다. 학점이나 과제 제출물의 차원을 넘어서서, 조사 시 배경 지식을 익히는 것도, 집중해서 알아보고자 하는 내용을 정하는 것도, 전사에 심혈을 기울이는 것도, 팀원과 협업하는 것도 모두 수강생의 몫입니다. 이 보고서의 이곳저곳이 장래의 수강생분들로 하여금 조금이나마 더 많은 것을 얻어 가시도록 도울 수 있다면 좋겠습니다.

\subsection{박준영}

언어조사를 준비하고, 실행하고, 자료를 정리하는 작업은 생각보다 고된 일입니다. 이전에 언어조사를 다녀온 많은 학우들이 그러했듯, 여러분도 헤드폰을 쓴 채 며칠 밤을 새우고, 같은 소리를 백 번씩 들으며 고뇌하고, 팀원들과 함께 머리를 싸매며 토론하는 시간을 보내게 될 것입니다. 때로는 언어조사 작업에 지치고, 수강취소를 하지 않은 과거의 자신을 원망하게 될지도 모릅니다.

그럼에도, 살아 있는 언어를 다루는 것은, 수(數)로 환산할 수 없을 정도로 대단히 값지고 의미 있는 일입니다. 매 순간 벌교 지역어의 마지막 모습을 남긴다는 사명감으로 언어조사에 임했고, 그 마음이 이 글을 완성하고 종강하는 데까지 저를 이끌었습니다. 미래에 `언어조사 및 분석'을 수강하실 학우 여러분께서도 사라져 가는 한국의 지역어를 담아내는 최전선에 서 있다는 자부심을 가지고, 힘 닿는 데까지 언어조사에 참여해 주세요. 우리의 보고서가 그 여정에 조금이나마 보탬이 되기를 진심으로 바랍니다.

\subsection{유승민}

저는 매 학기 수강신청 시기마다 듣고 싶은 수업들을 마음이 가는 대로 수강신청하는 편입니다. 언어조사 및 분석을 수강신청한 것도 그의 일환이었습니다.  애초에 현장언어학에 대한 선망이 있었고, 어찌 됐든, 결국에 품을 들이면 못 해낼 게 없다는 생각을 가지고 있었습니다. 이 강의를 수강하고자 하는 여러분 중에도 그러한 분들이 많으실 것이라고 생각합니다. 뭔가 재미있어 보이잖아요? 

\begin{quote}
    ``언어학과에는 수학여행을 가는 수업이 있대!''
\end{quote}

훨씬 로드가 적고, 몸과 마음이 편안한 다른 강의들이 (많이) 있습니다. 그럼에도 불구하고, 언어조사 및 분석을 수강신청하시거나 수강하시는 것에는, 현장언어학에 대한 강한 흥미가 기저에 있을 것입니다. 맞아요. 저도 정말 재미있었고, 이번 한 학기동안 이 강의에서 가장 많은 것을 얻어간 것 같습니다. 이 강의는 수강할 가치가 충분합니다. 아예 언어학과 커리큘럼의 중심에 놓고 싶을 정도입니다. 맞아요. 한번쯤은 들어보세요. 강력 추천합니다. 물귀신 작전이 아니에요. 

그러나, 여러분의 `로드를 버틸 수 있다는 배짱'에 조금의 태클을 걸고 싶네요. 이 과목은 제가 들은 어떠한 강의보다도 묵직했습니다. 섣부른 수강신청 이전에, 혹은 수강 포기 기간이 끝나기 이전에, 저는 여러분의 발목을 한번은 잡고 싶네요. 흥미가 가는 수업, 화기애애한 언어 조사 이후에 여러분은 보고서 작성을 시작하셔야 합니다. 그리고 여러분은 분명 보고서의 첫 문단을 전사하시면서 `무언가 단단히 잘못되었다는 것'을 느끼실 겁니다. 특히 저의 경우에는 바쁜 한 학기를 사느라, 좀처럼 보고서 작성에 전념할 시간이 나오지 않았습니다. 그런데요, 보고서 작성은 거대한 조별과제입니다. 팀플을 포기할 수는 없죠, 영원한 빌런이 되고 싶지 않다면 말입니다. 자연스럽게 보고서 작성은 여러분 의식 내에서 `1순위로 처리해야 할 과제'로 자리잡게 될 것입니다. 그러나, 보고서 작성은 쉽사리 끝나지 않습니다. 체감상 텀페이퍼를 3개는 쓰는 느낌입니다. 저처럼 힘든 학기에 이 수업을 수강신청하셨다면, 저와 같이 새벽에 전사를 하다가, 중간에 지쳐 쓰러지는 일을 겪으실 겁니다.

그렇기에 조언합니다. 바쁜 학기에는 수강신청하지 마세요. 과거의 제게 조언할 수 있다면, 저는 이번 학기에 6학점 정도를 덜 들으라고 할 것입니다. 그리고, 그저 학점을 채우겠다는 마음가짐보다는 현장언어학에 대한 선망과 애정을 가지고 오시기를 권합니다. 그저 팀플이라는 것 이외에도, 실제로 사라져가는 언어를 보존한다는 책임감도 뒤따르는 과목이니까요. 그러나, 저는 여러분께 한번쯤은 이 강의를 수강하시기를 강력하게 권합니다. 안락 의자에서 하지 않는 진짜 언어학을 할 수 있는 기회는 드무니까요.

p.s. 막판에 밤을 새고 싶지 않으시다면, 처음에 보고서 작성 계획을 과도할 정도로 세세하게 세우시기를 권합니다. (e.g. 3시간 단위의 목표 설정)

\subsection{이상인}

너희도 당해 봐라.

\subsection{이채현}

실제 공간에서 방언의 모습을 직접 관찰하고 방언의 화자분들과 상호작용하며 언어를 배워나가는 경험은 무척 값진 경험입니다. 다만 언어조사를 하다보면 후반부에 생각보다 시간 투자를 많이 해야 하고, 미리 조금씩 해두지 않으면 기말 기간에 전사 폭탄을 맞이하게 될 것입니다. 실제로 저희팀은 전사를 빠르게 마무리 하기 위해 전사지옥을 열고, 6시간 간격 중간 보고 시스템을 도입하는 등 노력했습니다. 가급적이면 한 학기에 듣는 수업 수가 적을 때 수강하는 것을 추천드립니다.

\subsection{전지민}

전사 계획을 짤 때 후반부로 갈 수록 기말고사 혹은 기타 학기말 일정과 겹쳐 굉장히 바빠지게 되기 때문에, 최대한 전반부에 많은 일을 할 수 있게 계획을 짜는 것이 좋을 것 같습니다. 또한 갑작스럽게 여러 화자분을 조사할 상황을 대비해 녹음 장비는 최대한 많이 준비해 가는 것이 좋을 것 같습니다.